% Generated by scripts/assemble.py; corrections belong in source/corrections.json.
\examhead{154}{中山大学 2026 年数学分析试卷}{146}
% p147-b16; source PDF page 147
\par\Needspace{5\baselineskip}
\blockanchor{p147-b16}
\majorquestion{一}填空题. 每题 5 分, 共 60 分.\par
% p147-b17; source PDF page 147
\blockanchor{p147-b17}
\question{1}求极限  $\lim_{x\to0^{+}}\left(\frac{\sin x}{x}\right)^{\frac{1}{x^{2}}}=\underline{\quad}$ .\par
% p147-b18; source PDF page 147
\blockanchor{p147-b18}
\question{2}设函数 $y = y(x)$ 由方程 $\ln(xy) = x^2 - y^2$ 隐式定义, 求点 $(x, y) = (1, 1)$ 处的导数 $\frac{\mathrm{d}y}{\mathrm{d}x}(1, 1) =$ \_\_\_\_.\par
% p147-b19; source PDF page 147
\blockanchor{p147-b19}
\question{3}求定积分  $\int_{0}^{\pi}\frac{dx}{\sqrt{2}+\sin x}=$  \_\_\_\_.\par
% p147-b20; source PDF page 147
\blockanchor{p147-b20}
\question{4}曲线由参数方程 $x = t, y = t^2$ 给出，求该曲线在 $t = 1$ 处的曲率\par
% p147-b21; source PDF page 147
\blockanchor{p147-b21}
\question{5}已知函数 $f(x, y) = \begin{cases} \frac{xy}{(x^2 + y^2)^p}, & x^2 + y^2 \neq 0 \\ 0, & x^2 + y^2 = 0 \end{cases}$ 在 (0,0) 处连续，求 $p$ 的取值范围 \_\_\_\_.\par
% p147-b22; source PDF page 147
\blockanchor{p147-b22}
\question{6}求函数 $\mu(x, y, z) = xy^2 + z^3 - xyz$ 在点 (1, 1, 2) 处沿方向 $l = \left(\frac{1}{2}, \frac{\sqrt{2}}{2}, \frac{1}{2}\right)$ 的方向导数 \_\_\_\_.\par
% p148-b00; source PDF page 148
\blockanchor{p148-b00}
\question{7}计算积分  $\iint_{0\leq s\leq1,s\leq t\leq\sqrt{s}}\frac{\sin t}{t}dtds=$  \_\_\_\_.\par
% p148-b01; source PDF page 148
\blockanchor{p148-b01}
\question{8}设 L 为球面  $x^{2} + y^{2} + z^{2} = 4$  被平面  $x + y + z = 0$  所截得的圆周, 计算积分  $\int_{L} x^{2} ds =$ \_\_\_\_ .\par
% p148-b02; source PDF page 148
\blockanchor{p148-b02}
\question{9}级数  $\sum_{n=2}^{\infty}\frac{1^{n}+2^{n}+\cdots+10^{n}}{n(\ln n)^{2}}e^{nx}$  的一致收敛区域为 \_\_\_\_.\par
% p148-b03; source PDF page 148
\blockanchor{p148-b03}
\question{10}如果级数 $\sum_{n=1}^{\infty}\left(1 - \frac{p\ln n}{2n}\right)^n$ 是收敛的, 则 $p$ 的取值范围是 \_\_\_\_.\par
% p148-b04; source PDF page 148
\blockanchor{p148-b04}
\question{11}设 $S(x) = \sum_{n=1}^{\infty} \frac{\sin 2nx}{n}$, 则 $S\left(\frac{\pi}{6}\right) = \underline{\quad}$.\par
% p148-b05; source PDF page 148
\blockanchor{p148-b05}
\question{12}$\int_{0}^{1}\frac{1}{x}\ln \frac{1 + x}{1 - x}\mathrm{d}x =$\par
% p148-b06; source PDF page 148
\par\Needspace{5\baselineskip}
\blockanchor{p148-b06}
\majorquestion{二}(15 分) 设 $f(x)$ 在 (0,1) 上有定义, $f(x) \tan x$ 单调递增, $\tan(f(x))$ 单调递减, 证明: $f(x)$ 在 (0,1) 上连续.\par
% p148-b07; source PDF page 148
\par\Needspace{5\baselineskip}
\blockanchor{p148-b07}
\majorquestion{三}(15 分) 证明: 函数 $f(x) = x \cos \frac{1}{x}$ 在 $(0, +\infty)$ 上一致连续.\par
% p148-b08; source PDF page 148
\par\Needspace{5\baselineskip}
\blockanchor{p148-b08}
\majorquestion{四}(15 分) 计算积分 $I(b) = \int_{0}^{\frac{\pi}{2}} \frac{\mathrm{d}x}{b + \tan x}, b > 0$.\par
% p148-b09; source PDF page 148
\par\Needspace{5\baselineskip}
\blockanchor{p148-b09}
\majorquestion{五}(15 分) 设 $f(x)$ 在 $[0, +\infty)$ 上一致连续, 且对任意固定的 $x \in [0,1]$, 有 $\lim_{n \to \infty} f(x + n) = 0$, 证明: 函数列 $\{f(x + n) \mid n \geq 1\}$ 在 $[0,1]$ 上一致收敛于 0.\par
% p148-b10; source PDF page 148
\par\Needspace{5\baselineskip}
\par\Needspace{6\baselineskip}
\blockanchor{p148-b10}
\majorquestion{六}(15 分) 设函数列 $\{a_{n}(x)\}$ 在 $[0,1]$ 上有定义, 且存在不依赖于 $x, n$ 的常数 $K$, 使得\par
\nopagebreak[4]
% p148-b11; source PDF page 148
\blockanchor{p148-b11}
\displayformula{\left| a _ {1} (x) \right| \leq K, \sum_ {n = 1} ^ {m} \left| a _ {n} (x) - a _ {n + 1} (x) \right| \leq K, \forall m \geq 1.}
% p148-b12; source PDF page 148
\blockanchor{p148-b12}
证明: 如果级数  $\sum_{n=1}^{\infty}c_{n}$  收敛,  $c_{n}\geq0$ , 则级数  $\sum_{n=1}^{\infty}c_{n}a_{n}(x)$  在  $[0,1]$  上一致收敛.\par
% p148-b13; source PDF page 148
\par\Needspace{5\baselineskip}
\blockanchor{p148-b13}
\majorquestion{七}(15 分) 令 $B_{R} = \{(x_{1}, x_{2}) \in \mathbb{R}^{2} \mid x_{1}^{2} + x_{2}^{2} \leq R^{2}\}$ 是 $\mathbb{R}^{2}$ 中的半径为 $R$ 的闭球, $R > 0$. 设 $f(x)$ 是 $B_{2}$ 上的连续函数, 满足条件: 对任意的 $x, y \in B_{2}$, 成立 $\frac{f(x) + f(y)}{2} \geq f\left(\frac{x + y}{2}\right)$. 证明:\par
% p148-b14; source PDF page 148
\blockanchor{p148-b14}
(1) 对任意的 $x, y \in B_2$ 和任意的 $\lambda \in (0, 1)$, 成立 $\lambda f(x) + (1 - \lambda)f(y) \geq f(\lambda x + (1 - \lambda)y)$.\par
% p148-b15; source PDF page 148
\par\Needspace{6\baselineskip}
\blockanchor{p148-b15}
(2) 存在一个常数 $C > 0$, 使得如下不等式成立:\par
\nopagebreak[4]
% p148-b16; source PDF page 148
\blockanchor{p148-b16}
\displayformula{| f (\boldsymbol {x}) - f (\boldsymbol {y}) | \leq C \| \boldsymbol {x} - \boldsymbol {y} \|, \forall \boldsymbol {x}, \boldsymbol {y} \in B _ {1}.}
% p148-b17; source PDF page 148
\blockanchor{p148-b17}
这里 $\| \pmb {x} - \pmb{y}\|$ 表示 $\pmb {x} = (x_1,x_2)$ 和 $\pmb {y} = (y_{1},y_{2})$ 之间的距离 $\sqrt{(x_1 - y_1)^2 + (x_2 - y_2)^2}.$\par
